\scanpage{113}
alors qu'il est clair que $\inserted{\unsure{z_0}=}x_0$, $\inserted{z_0=}y_0$ sont solutions solutions de l'équation (48). Si donc $z_0$ est une solution $\ne x_0,y_0$, on aura
\[
 \begin{aligned}
 x_0+y_0+z_0&=-(a_1-\lambda^2)=\lambda^2-a_1\\
 &=\left(\frac{x'-y'}{x_0-y_0}\right)^2\unsure{-|=}a_1
 \end{aligned}
\]
d'où la solution (sous réserve d'existence d'une solution $z_0\ne x_0,y_0$, ou encore \deleted{\unsure{d'une so|d'autre part}} \inserted{sous} réserve que
\[
 \left\{\begin{aligned}
 &x,y\in E\smallsetminus\{0\},\quad x\ne y,\quad -(x+y)\ne x,y\quad\text{i.e.}\\
 &\hspace{10em}2x+y\ne0,\quad 2y+x\ne0\\
 &\hspace{10em}\text{i.e.}\quad y\ne-2x,\ x\ne-2y\quad),
 \end{aligned}\right. \tag*{(50)}
\]
la solution
\[
 \left\{\begin{aligned}
 z_0&=-x_0-y_0+\left(\frac{x'-y'}{x_0-y_0}\right)^2-a_1\\
 z'&=\lambda z_0+\mu=\frac{x'-y'}{x_0-y_0}
 \left[-x_0-y_0+\left(\frac{x'-y'}{x_0-y_0}\right)^2-a_1\right]\\
 &\hspace{9em}\inserted{+\frac{y'x_0-x'y_0}{x_0-y_0}}\,.
 \end{aligned}\right. \tag*{(51)}
\]
Par principe de prolongement, cette expression pour $z_0=p(z)$ et $z'=p'(z)$, pour $z=-(x\unsure{-|+}y)$, est valable encore dès que $x,y\in E(k)\smallsetminus\{0\}$ sont tels que $x_0\ne y_0$ i.e. $p(x)\ne p(y)$, i.e. dès que $\deleted{x\unsure{=|\ne}}y\notin\pm x$, de sorte que les formules (51) ont un sens.
% archive sheet number: 112
